% !TEX program = pdflatex
% THEME 2 — Moyen — Exercices
\documentclass[11pt,a4paper]{article}
\usepackage[utf8]{inputenc}
\usepackage[T1]{fontenc}
\usepackage{lmodern}
\usepackage[french]{babel}
\usepackage{amsmath,amssymb,mathtools}
\usepackage{icomma}
\usepackage{siunitx}
\sisetup{output-decimal-marker={,},per-mode=symbol}
\DeclareSIUnit{\atm}{atm}
\DeclareSIUnit{\bar}{bar}
\DeclareSIUnit{\mmHg}{mmHg}
\usepackage[version=4]{mhchem}
\usepackage{xcolor}
\usepackage{colortbl}
\usepackage{tikz}
\usepackage[most]{tcolorbox}
\usepackage{enumitem}
\usepackage{array}
\usepackage{booktabs}
\usepackage{fancyhdr}
\usepackage[a4paper,margin=2.1cm,headheight=26pt,headsep=10pt]{geometry}
\usetikzlibrary{arrows.meta,calc,shapes.geometric}

\definecolor{primary}{RGB}{150,98,162}
\definecolor{primarydark}{RGB}{92,52,110}
\definecolor{gazcol}{RGB}{0,120,190}
\definecolor{azote}{RGB}{60,90,200}
\definecolor{oxygene}{RGB}{210,60,60}

\newtcolorbox{enonce}[1][]{enhanced,breakable,colback=primary!4,colframe=primary,
  boxrule=0.8pt,arc=2pt,left=3mm,right=3mm,top=2mm,bottom=2mm,
  fonttitle=\bfseries,coltitle=white,title={#1}}

\pagestyle{fancy}\fancyhf{}
\fancyhead[L]{\small\textcolor{primarydark}{\textbf{Fiche 7} — Loi des gaz}}
\fancyhead[R]{\small\textcolor{primarydark}{\textsc{Thème 2 · Moyen · Exercices}}}
\fancyfoot[C]{\small\thepage}
\renewcommand{\headrulewidth}{0.8pt}
\renewcommand{\headrule}{\hbox to\headwidth{\color{primary}\leaders\hrule height \headrulewidth\hfill}}
\setlength{\parindent}{0pt}\setlength{\parskip}{3pt}

\newcounter{exo}
\newenvironment{exo}[1][]{\refstepcounter{exo}\medskip
  \noindent{\color{primarydark}\bfseries Exercice \theexo.}\ \textit{#1}\par\nopagebreak\smallskip}{\par}

\begin{document}

\begin{center}
{\Large\bfseries\color{primarydark} Thème 2 — Niveau Moyen}\\[2pt]
{\color{primary}\itshape Passer par la masse, comparer deux enceintes, surveiller les unités}
\end{center}
\textcolor{primary}{\hrule height 1pt}
\medskip

\begin{enonce}[Données]
$R=\qty{0.082}{\atm\litre\per\mole\per\kelvin}$ ; \ $n=\dfrac{m}{M}$ ; \
$M(\ce{O2})=\qty{32}{\gram\per\mole}$, \ $M(\ce{CH4})=\qty{16}{\gram\per\mole}$,
\ $M(\ce{N2})=\qty{28}{\gram\per\mole}$.
\end{enonce}

\begin{exo}[Peser un gaz]
Une bouteille de plongée rigide de volume $V=\qty{10.0}{\litre}$ contient du dioxygène \ce{O2} sous une
pression $p=\qty{5.0}{\atm}$ à $\theta=\qty{27}{\degreeCelsius}$.
\begin{enumerate}[label=(\alph*),leftmargin=2em,itemsep=1pt]
  \item Calcule la quantité de matière $n$ de dioxygène.
  \item Quelle masse $m$ de dioxygène la bouteille contient-elle ?
\end{enumerate}
\end{exo}

\begin{exo}[Retrouver une température]
Un récipient rigide de volume $V=\qty{2.0}{\litre}$ contient $n=\qty{0.10}{\mole}$ de gaz sous une
pression $p=\qty{1.5}{\atm}$. Quelle est la température $T$ du gaz (en kelvins, puis en degrés Celsius) ?
\end{exo}

\begin{exo}[L'unité qui change tout]
Un gaz occupe $V=\qty{5.0}{\litre}$ sous $p=\qty{1.0}{\atm}$ à $T=\qty{300}{\kelvin}$. Un étudiant calcule
la quantité de matière avec $R=\num{8.314}$ et trouve $n\approx\qty{0.0020}{\mole}$.
\begin{enumerate}[label=(\alph*),leftmargin=2em,itemsep=1pt]
  \item Pourquoi ce résultat est-il faux ? Quel $R$ fallait-il prendre ?
  \item Donne la valeur correcte de $n$, et le facteur d'erreur commis par l'étudiant.
\end{enumerate}
\end{exo}

\begin{exo}[Deux enceintes à la même température]
Du méthane gazeux \ce{CH4} de masse $\qty{32}{\gram}$ est enfermé dans une enceinte de volume
$\qty{5.0}{\litre}$ sous une pression de $\qty{0.20}{\atm}$. Du dioxygène gazeux \ce{O2} de masse
$\qty{32}{\gram}$ est enfermé dans une autre enceinte de volume $\qty{10.0}{\litre}$.

\begin{center}
\begin{tikzpicture}[font=\footnotesize]
  \draw[thick] (0,0) rectangle (2.2,2.2); \fill[gazcol!8](0,0)rectangle(2.2,2.2);
  \foreach \x/\y in {0.5/0.5,1.2/1.4,1.7/0.7,0.8/1.8}{\fill[azote](\x,\y)circle(2.6pt);}
  \node[below] at (1.1,-0.15){\ce{CH4} : $\qty{32}{\gram}$};
  \node[below] at (1.1,-0.6){$\qty{5.0}{\litre}$, $\qty{0.20}{\atm}$};
  \node[thick] at (3.3,1.1){\Large ?};
  \draw[thick] (4.4,0) rectangle (6.6,2.2); \fill[gazcol!8](4.4,0)rectangle(6.6,2.2);
  \foreach \x/\y in {4.9/0.6,5.6/1.5,6.1/0.7,5.1/1.9,6.2/1.6}{\fill[oxygene](\x,\y)circle(3pt);}
  \node[below] at (5.5,-0.15){\ce{O2} : $\qty{32}{\gram}$};
  \node[below] at (5.5,-0.6){$\qty{10.0}{\litre}$, $p=?$};
\end{tikzpicture}
\end{center}
\vspace{0.5cm}

Quelle doit être la pression $p$ dans l'enceinte contenant le dioxygène pour que la \textbf{température
soit la même} dans les deux enceintes ? \emph{(Indication : à même $T$, la quantité $pV/n$ est identique
pour les deux gaz ; inutile de calculer $T$.)}
\end{exo}

\begin{exo}[Le corollaire d'Avogadro, sans nombre]
Deux ballons \textbf{identiques} (même volume) contiennent, à la \textbf{même température} et sous la
\textbf{même pression}, l'un du dihydrogène \ce{H2}, l'autre du dioxyde de carbone \ce{CO2}. En partant de
$pV=nRT$, montre qu'ils contiennent le \emph{même nombre de moles}, et donc le même nombre de molécules,
bien que les masses de gaz soient très différentes.
\end{exo}

\end{document}
